% GENERATED FILE. Edit canonical lessons, then run npm run generate:books.
% canonical-source-hash: fnv1a64:116b826d22ab87db
% canonical-lessons: RU-C33-temperature, RU-C33-cough, RU-C33-rest, RU-C33-hunger, RU-C33-thirst

\chapter{Fever, Cough, Rest, Hunger, Thirst}
\label{ch:ru-fever-cough-rest-hunger-thirst}
\hlchaptermodality{33}

\begin{chapteropening}
\textbf{By the end of this chapter:} \emph{I can name a temperature, a cough, rest, hunger and thirst.}

Complete the last lesson of chapter 33: I can name a temperature, a cough, rest, hunger and thirst.
\end{chapteropening}

\section[\texorpdfstring{\ru{температура} (temperatúra)}{температура (temperatúra)}]{\ru{температура} (temperatúra) — a temperature, a fever}

\label{lesson:RU-C33-temperature}



\begin{quote}
Before the new one: say the Russian for a doctor, then the Russian for a hospital.
\end{quote}

\subsection*{You'll want to know: \ru{температура}}
\textbf{\ru{температура}} (\emph{temperatúra}) — \textquotedblleft{}a temperature, a fever\textquotedblright{}.

\begin{cousinweb}
\textbf{\ru{температура}} $\leftarrow$ Latin \emph{temperātūra}, \textquotedblleft{}due measure, mixing\textquotedblright{}, from \emph{temperāre}, \textquotedblleft{}to mix in proportion\textquotedblright{} — borrowed in modern times, probably through German \emph{Temperatur}.
\end{cousinweb}

\begin{grammarlens}[title={What you've built}]
One more word for this chapter.
\end{grammarlens}

\subsection*{Your turn}
\begin{itemize}
\raggedright
  \item \emph{Say it:} \emph{temperatúra}
  \item \emph{Say it:} \emph{temperatúra}, once more
  \item \emph{Say it:} say \emph{temperatúra}
  \item \emph{Recall:} say the Russian for a doctor, then the Russian for a hospital, then say \emph{temperatúra} again
\end{itemize}

\begin{tcolorbox}[breakable,colback=teal!4,colframe=teal!35!black,title={Before you move on}]
What does \ru{температура} mean? (\textquotedblleft{}A temperature, a fever\textquotedblright{}.) Say it once more.
\end{tcolorbox}
\section[\texorpdfstring{\ru{кашель} (káshel)}{кашель (káshel)}]{\ru{кашель} (káshel) — a cough}

\label{lesson:RU-C33-cough}



\begin{quote}
Before the new one: say the Russian for a hospital, then the Russian for a temperature.
\end{quote}

\subsection*{You'll want to know: \ru{кашель}}
\textbf{\ru{кашель}} (\emph{káshel}) — \textquotedblleft{}a cough\textquotedblright{}.

\begin{cousinweb}
\textbf{\ru{кашель}} $\leftarrow$ Proto-Slavic \emph{*kašĭlĭ} — inherited, a cousin of Lithuanian \emph{kosulỹs}, Sanskrit \emph{kāsa-} and Old High German \emph{huosto}, all \textquotedblleft{}cough\textquotedblright{} (Vasmer).
\end{cousinweb}

\begin{grammarlens}[title={What you've built}]
One more word for this chapter.
\end{grammarlens}

\subsection*{Your turn}
\begin{itemize}
\raggedright
  \item \emph{Say it:} \emph{káshel}
  \item \emph{Say it:} \emph{káshel}, once more
  \item \emph{Say it:} say \emph{káshel}
  \item \emph{Recall:} say the Russian for a hospital, then the Russian for a temperature, then say \emph{káshel} again
\end{itemize}

\begin{tcolorbox}[breakable,colback=teal!4,colframe=teal!35!black,title={Before you move on}]
What does \ru{кашель} mean? (\textquotedblleft{}A cough\textquotedblright{}.) Say it once more.
\end{tcolorbox}
\section[\texorpdfstring{\ru{отдых} (ótdykh)}{отдых (ótdykh)}]{\ru{отдых} (ótdykh) — rest}

\label{lesson:RU-C33-rest}



\begin{quote}
Before the new one: say the Russian for a temperature, then the Russian for a cough.
\end{quote}

\subsection*{You'll want to know: \ru{отдых}}
\textbf{\ru{отдых}} (\emph{ótdykh}) — \textquotedblleft{}rest\textquotedblright{}.

\begin{cousinweb}
\textbf{\ru{отдых}} is \emph{\ru{от}-}, \textquotedblleft{}back, away\textquotedblright{}, with \emph{\ru{дых}}, \textquotedblleft{}breath\textquotedblright{}, from the verb \emph{\ru{дышать}}, \textquotedblleft{}to breathe\textquotedblright{}: rest is getting your breath back. For the root, Vasmer compares Lithuanian \emph{dūsėti}, \textquotedblleft{}to pant\textquotedblright{}.
\end{cousinweb}

\begin{grammarlens}[title={What you've built}]
One more word for this chapter.
\end{grammarlens}

\subsection*{Your turn}
\begin{itemize}
\raggedright
  \item \emph{Say it:} \emph{ótdykh}
  \item \emph{Say it:} \emph{ótdykh}, once more
  \item \emph{Say it:} say \emph{ótdykh}
  \item \emph{Recall:} say the Russian for a temperature, then the Russian for a cough, then say \emph{ótdykh} again
\end{itemize}

\begin{tcolorbox}[breakable,colback=teal!4,colframe=teal!35!black,title={Before you move on}]
What does \ru{отдых} mean? (\textquotedblleft{}Rest\textquotedblright{}.) Say it once more.
\end{tcolorbox}
\section[\texorpdfstring{\ru{голод} (gólod)}{голод (gólod)}]{\ru{голод} (gólod) — hunger}

\label{lesson:RU-C33-hunger}



\begin{quote}
Before the new one: say the Russian for a cough, then the Russian for rest.
\end{quote}

\subsection*{You'll want to know: \ru{голод}}
\textbf{\ru{голод}} (\emph{gólod}) — \textquotedblleft{}hunger\textquotedblright{}.

\begin{cousinweb}
\textbf{\ru{голод}} $\leftarrow$ Proto-Slavic \emph{*goldŭ} — inherited. Vasmer links it with Church Slavonic \emph{žĭlděti}, \textquotedblleft{}to crave\textquotedblright{}, and Sanskrit \emph{gṛdhyati}, \textquotedblleft{}is greedy\textquotedblright{}. The Church Slavonic shape is \emph{\ru{глад}}, \textquotedblleft{}famine\textquotedblright{}, with \emph{-\ru{ла}-} for Russian \emph{-\ru{оло}-}.
\end{cousinweb}

\begin{grammarlens}[title={What you've built}]
One more word for this chapter.
\end{grammarlens}

\subsection*{Your turn}
\begin{itemize}
\raggedright
  \item \emph{Say it:} \emph{gólod}
  \item \emph{Say it:} \emph{gólod}, once more
  \item \emph{Say it:} say \emph{gólod}
  \item \emph{Recall:} say the Russian for a cough, then the Russian for rest, then say \emph{gólod} again
\end{itemize}

\begin{tcolorbox}[breakable,colback=teal!4,colframe=teal!35!black,title={Before you move on}]
What does \ru{голод} mean? (\textquotedblleft{}Hunger\textquotedblright{}.) Say it once more.
\end{tcolorbox}
\section[\texorpdfstring{\ru{жажда} (zházhda)}{жажда (zházhda)}]{\ru{жажда} (zházhda) — thirst}

\label{lesson:RU-C33-thirst}



\begin{quote}
Before the new one: say the Russian for rest, then the Russian for hunger.
\end{quote}

\subsection*{You'll want to know: \ru{жажда}}
\textbf{\ru{жажда}} (\emph{zházhda}) — \textquotedblleft{}thirst\textquotedblright{}. \textbf{\ru{вода}} answers it.

\begin{cousinweb}
\textbf{\ru{жажда}} — taken from Church Slavonic, says Vasmer: \emph{-\ru{жд}-} is the Church Slavonic sound where East Slavic has \emph{-\ru{ж}-}, and Old Russian had the native \emph{\ru{жажа}}. Both go back to Proto-Slavic \emph{*žędja}.
\end{cousinweb}

\begin{grammarlens}[title={What you've built}]
One more word for this chapter.
\end{grammarlens}

\subsection*{Your turn}
\begin{itemize}
\raggedright
  \item \emph{Say it:} \emph{zházhda}
  \item \emph{Say it:} \emph{zházhda}, once more
  \item \emph{Say it:} say \emph{zházhda}
  \item \emph{Recall:} say the Russian for rest, then the Russian for hunger, then say \emph{zházhda} again
\end{itemize}

\begin{tcolorbox}[breakable,colback=teal!4,colframe=teal!35!black,title={Before you move on}]
What does \ru{жажда} mean? (\textquotedblleft{}Thirst\textquotedblright{}.) Say it once more.
\end{tcolorbox}
